\documentclass[11pt,a4paper]{article}
\usepackage[utf8]{inputenc}
\usepackage[T1]{fontenc}
\usepackage[portuguese]{babel}
\usepackage{xcolor}
\usepackage{hyperref}
\usepackage{geometry}
\usepackage{tcolorbox}

\geometry{margin=2.5cm}

\definecolor{primary}{RGB}{39, 174, 96}
\definecolor{secondary}{RGB}{44, 62, 80}
\definecolor{codebg}{RGB}{240, 240, 240}

\hypersetup{colorlinks=true, linkcolor=primary, urlcolor=primary}

\title{\color{secondary}\Huge\textbf{Solução: Exercício 4.6 - O projeto, passo 23}}
\author{\color{primary}DevOps com Kubernetes -- 2026}
\date{}

\begin{document}
\maketitle

\section*{\color{primary}Guia de Solução}
\begin{tcolorbox}[colback=green!5!white,colframe=green!60!black,title=Resolução Passo a Passo]
### Passo a Passo

1. \textbf{Publicar Mensagens (Backend)}: Integre a biblioteca do NATS no backend principal. Logo após inserir/atualizar um Todo no banco de dados com sucesso, utilize \texttt{nc.publish('todo_events', JSON.stringify(todoData))}.
2. \textbf{Criar o Serviço Broadcaster}: Crie um novo microsserviço (ex: Node, Python, Go) do zero. Este serviço conectará ao NATS e fará a inscrição (\textit{subscribe}) no tópico \texttt{todo_events}.
3. \textbf{Encaminhar para Serviço Externo}: Quando o \textit{broadcaster} receber uma mensagem da fila, ele deverá fazer uma requisição POST HTTP (ex: usando \texttt{axios} ou \texttt{curl}) para a URL de um \textit{webhook} do Discord/Telegram (ou um \textit{webhook} de testes genérico) com os dados formatados.
4. \textbf{Implantação no K8s}: Crie um \texttt{Deployment} para o \textit{broadcaster} garantindo que ele tenha a URL do NATS e a URL do webhook injetadas via variáveis de ambiente.
\end{tcolorbox}

\end{document}
